\subsection{Caution-aware model against the proposed model and both baselines (200 episodes, four weather conditions)}\label{sec:cautioncompare}
Same protocol as Section~\ref{sec:headline}: scenario seeds 5000 to 5199, sparse traffic, 800 steps, weather applied as removal of vehicles from the grid at the measured detection rate. The caution-aware model receives its caution number from the weather (clear weather: caution 0, which is identical to its clear-weather behaviour). The time limit was \emph{not} extended; a timeout counts as a failure.
\begin{table}[H]\centering\scriptsize
\caption{Every metric for every model and weather condition, 200 episodes each. Collision, left road and timeout are shares of episodes; the success rate is one minus their sum. Energy is an estimate for a declared virtual vehicle without regenerative braking.}\label{tab:cautioncompare}
\resizebox{\textwidth}{!}{\begin{tabular}{@{}llrrrrrrrr@{}}\toprule
Weather & Model & Success & 95\% interval & Collision & Left road & Timeout & Speed (m/s) & Route completion & Energy (Wh per km)\\\midrule
Clear & Baseline Driving Model Trained Without Traffic & 0.850 & [0.794, 0.893] & 0.125 & 0.025 & 0.000 & 8.09 & 0.918 & 85.9\\
 & Baseline Driving Model Trained With Traffic & 0.745 & [0.680, 0.800] & 0.120 & 0.135 & 0.000 & 8.07 & 0.883 & 90.7\\
 & Proposed Bird's-Eye-View Driving Model & 0.925 & [0.880, 0.954] & 0.060 & 0.010 & 0.005 & 7.87 & 0.939 & 96.2\\
 & Caution-Study Caution-Aware Driving Model (caution from weather) & 0.890 & [0.839, 0.926] & 0.090 & 0.020 & 0.000 & 7.97 & 0.920 & 94.1\\
\midrule
Fog (intensity 1.0) & Baseline Driving Model Trained Without Traffic & 0.860 & [0.805, 0.901] & 0.115 & 0.025 & 0.000 & 8.09 & 0.920 & 85.9\\
 & Baseline Driving Model Trained With Traffic & 0.750 & [0.686, 0.805] & 0.120 & 0.130 & 0.000 & 8.07 & 0.886 & 90.5\\
 & Proposed Bird's-Eye-View Driving Model & 0.895 & [0.845, 0.930] & 0.085 & 0.020 & 0.000 & 7.92 & 0.933 & 96.5\\
 & Caution-Study Caution-Aware Driving Model (caution from weather) & 0.895 & [0.845, 0.930] & 0.050 & 0.020 & 0.035 & 6.15 & 0.945 & 62.7\\
\midrule
Rain (intensity 0.6) & Baseline Driving Model Trained Without Traffic & 0.850 & [0.794, 0.893] & 0.125 & 0.025 & 0.000 & 8.09 & 0.918 & 86.0\\
 & Baseline Driving Model Trained With Traffic & 0.755 & [0.691, 0.809] & 0.120 & 0.125 & 0.000 & 8.08 & 0.886 & 90.5\\
 & Proposed Bird's-Eye-View Driving Model & 0.900 & [0.851, 0.934] & 0.085 & 0.015 & 0.000 & 7.94 & 0.934 & 94.9\\
 & Caution-Study Caution-Aware Driving Model (caution from weather) & 0.890 & [0.839, 0.926] & 0.030 & 0.020 & 0.060 & 5.87 & 0.951 & 60.4\\
\midrule
Rain (intensity 1.0) & Baseline Driving Model Trained Without Traffic & 0.850 & [0.794, 0.893] & 0.125 & 0.025 & 0.000 & 8.09 & 0.918 & 85.9\\
 & Baseline Driving Model Trained With Traffic & 0.755 & [0.691, 0.809] & 0.120 & 0.125 & 0.000 & 8.08 & 0.886 & 90.5\\
 & Proposed Bird's-Eye-View Driving Model & 0.900 & [0.851, 0.934] & 0.095 & 0.005 & 0.000 & 7.95 & 0.932 & 93.0\\
 & Caution-Study Caution-Aware Driving Model (caution from weather) & 0.820 & [0.761, 0.867] & 0.035 & 0.040 & 0.105 & 5.56 & 0.935 & 59.7\\
\bottomrule\end{tabular}
}
\end{table}
\explainedfigure{fig_caution_comparison_200_episodes}{0.9}{Success, estimated energy and speed of the caution-aware model, the proposed model and the two baseline models in four weather conditions}{Each of the 200 test episodes (scenario seeds 5000 to 5199, sparse traffic) was driven by four models in clear weather, fog 1.0, rain 0.6 and rain 1.0, where the weather removes vehicles from the model's grid at the measured detection rate. Top: success rate with 95\% intervals. Middle: estimated traction energy in watt-hours per kilometre. Bottom: mean speed in metres per second.}{The caution-aware model reaches the same success as the proposed model in clear weather, fog and rain 0.6 (differences within $\pm0.035$, not significant) and is lower in rain 1.0 (0.820 against 0.900, $p=0.033$). In fog and rain it drives 22 to 30\% slower and uses 35 to 37\% less estimated energy, with fewer collisions and more timeouts. The two baselines hardly change between conditions, although their grid is thinned in the same way.}{Learned adaptive behaviour is real and saves energy (about 60 against 94 to 96 watt-hours per kilometre) without a significant loss of success in fog and moderate rain, but at the strongest rain it costs 0.080 of success through timeouts. The goal of not lowering the success rate is therefore met in clear weather, fog and rain 0.6, and not met in rain 1.0 within the 800-step limit.}{fig:cautioncompare}

\begin{table}[H]\centering\scriptsize
\caption{Paired comparison of the caution-aware model with each other model on the same 200 episodes. Positive differences favour the caution-aware model.}\label{tab:cautionpaired}
\resizebox{\textwidth}{!}{\begin{tabular}{@{}llrrrr@{}}\toprule
Weather & Compared model & Difference in success & Only caution-aware model succeeded & Only compared model succeeded & Exact McNemar $p$\\\midrule
Clear & Baseline Driving Model Trained Without Traffic & +0.040 & 11 & 3 & 0.0574\\
 & Baseline Driving Model Trained With Traffic & +0.145 & 32 & 3 & $<0.0001$\\
 & Proposed Bird's-Eye-View Driving Model & -0.035 & 3 & 10 & 0.0923\\
\midrule
Fog (intensity 1.0) & Baseline Driving Model Trained Without Traffic & +0.035 & 24 & 17 & 0.3489\\
 & Baseline Driving Model Trained With Traffic & +0.145 & 44 & 15 & 0.0002\\
 & Proposed Bird's-Eye-View Driving Model & +0.000 & 19 & 19 & 1.0000\\
\midrule
Rain (intensity 0.6) & Baseline Driving Model Trained Without Traffic & +0.040 & 29 & 21 & 0.3222\\
 & Baseline Driving Model Trained With Traffic & +0.135 & 45 & 18 & 0.0009\\
 & Proposed Bird's-Eye-View Driving Model & -0.010 & 20 & 22 & 0.8776\\
\midrule
Rain (intensity 1.0) & Baseline Driving Model Trained Without Traffic & -0.030 & 24 & 30 & 0.4966\\
 & Baseline Driving Model Trained With Traffic & +0.065 & 41 & 28 & 0.1480\\
 & Proposed Bird's-Eye-View Driving Model & -0.080 & 17 & 33 & 0.0328\\
\bottomrule\end{tabular}
}
\end{table}
\paragraph{Longer time limit (post-hoc check).} Because every extra failure of the caution-aware model was a timeout, it was re-evaluated in fog 1.0, rain 0.6 and rain 1.0 with a limit of 1600 steps (160 seconds) instead of 800 (same seeds, same dropout stream, \pth{bev_rl/adaptive/queue_caution_longer_time_limit.sh}). The proposed model had no timeouts in these conditions, so its result is unchanged by construction; clear weather had no caution-aware timeouts either. The limit was changed \emph{after} seeing the timeouts, so this is a post-hoc analysis, not a pre-registered one.
\begin{table}[H]\centering\scriptsize
\caption{Caution-aware model with the 800-step and the 1600-step limit, next to the proposed model (800 steps). 200 episodes each.}\label{tab:cautionlonger}
\resizebox{\textwidth}{!}{\begin{tabular}{@{}lllrrrrrr@{}}\toprule
Weather & Model & Step limit & Success & Collision & Left road & Timeout & Speed (m/s) & Energy (Wh per km)\\\midrule
Fog (intensity 1.0) & Proposed Bird's-Eye-View Driving Model & 800 (80 s) & 0.895 & 0.085 & 0.020 & 0.000 & 7.92 & 96.5\\
 & Caution-Study Caution-Aware Driving Model & 800 (80 s) & 0.895 & 0.050 & 0.020 & 0.035 & 6.15 & 62.7\\
 & Caution-Study Caution-Aware Driving Model & 1600 (160 s) & 0.940 & 0.050 & 0.010 & 0.000 & 6.15 & 62.5\\
\midrule
Rain (intensity 0.6) & Proposed Bird's-Eye-View Driving Model & 800 (80 s) & 0.900 & 0.085 & 0.015 & 0.000 & 7.94 & 94.9\\
 & Caution-Study Caution-Aware Driving Model & 800 (80 s) & 0.890 & 0.030 & 0.020 & 0.060 & 5.87 & 60.4\\
 & Caution-Study Caution-Aware Driving Model & 1600 (160 s) & 0.955 & 0.025 & 0.020 & 0.000 & 5.87 & 60.2\\
\midrule
Rain (intensity 1.0) & Proposed Bird's-Eye-View Driving Model & 800 (80 s) & 0.900 & 0.095 & 0.005 & 0.000 & 7.95 & 93.0\\
 & Caution-Study Caution-Aware Driving Model & 800 (80 s) & 0.820 & 0.035 & 0.040 & 0.105 & 5.56 & 59.7\\
 & Caution-Study Caution-Aware Driving Model & 1600 (160 s) & 0.935 & 0.035 & 0.030 & 0.000 & 5.56 & 59.7\\
\bottomrule\end{tabular}
}
\end{table}
\begin{table}[H]\centering\scriptsize
\caption{Paired comparison at the longer limit.}\label{tab:cautionlongerpaired}
\resizebox{\textwidth}{!}{\begin{tabular}{@{}lrrrr@{}}\toprule
Weather & Difference in success (caution-aware at 1600 minus proposed at 800) & Only caution-aware succeeded & Only proposed succeeded & Exact McNemar $p$\\\midrule
Fog (intensity 1.0) & +0.045 & 19 & 10 & 0.1360\\
Rain (intensity 0.6) & +0.055 & 20 & 9 & 0.0614\\
Rain (intensity 1.0) & +0.035 & 20 & 13 & 0.2962\\
\bottomrule\end{tabular}
}
\end{table}
With 160 seconds all timeouts disappear and the caution-aware model finishes 0.940, 0.955 and 0.935 of the episodes in fog 1.0, rain 0.6 and rain 1.0, against 0.895, 0.900 and 0.900 for the proposed model. None of the differences is significant ($p=0.136$, $0.061$, $0.296$), so the result is a tie or a small advantage, not a loss. The caution-aware model reaches the goal more slowly (the same speed and energy as before) but with fewer collisions (2.5 to 5\% against 8.5 to 9.5\%).

\begin{implicationbox}[What the comparison shows]
\textbf{Success.} The caution-aware model is statistically tied with the proposed model in clear weather ($-0.035$), fog ($0.000$) and rain 0.6 ($-0.010$), and significantly lower in rain 1.0 ($-0.080$, $p=0.033$). It beats the baseline trained with traffic in clear, fog and rain 0.6, and is not distinguishable from the baseline trained without traffic.
\textbf{Energy and speed.} In weather the caution-aware model uses 59.7 to 62.7 watt-hours per kilometre against 93.0 to 96.5 for the proposed model, about 35\% less, at 5.6 to 6.2 against 7.9 metres per second.
\textbf{Why success drops in rain 1.0.} All of the extra failures are timeouts (10.5\%): the slow driver does not finish within 80 seconds. Collisions fall from 9.5\% to 3.5\%. With a 160-second limit these timeouts vanish and rain 1.0 success becomes 0.935 (post-hoc, see the paragraph above), a tie with the proposed model.
\textbf{Why the two differ in clear weather.} The caution-aware model is not the proposed model with an extra input switched off: it is a different network, trained in a separate run on labels from a teacher whose speed and gap depend on caution, so even with caution 0 its weights and its driving differ. The proposed model succeeds in 185 of 200 clear episodes and the caution-aware model in 178; they disagree in 13 episodes (10 in favour of the proposed model, 3 against; $p=0.092$). The control model, trained like the caution-aware model but with caution fixed at 0, scores 0.895, the same as the caution-aware model. Separate training runs of the same recipe therefore differ by a few percentage points at 200 episodes, and 0.925 is at the favourable end of that spread; with one training run per model this cannot be resolved further.
\textbf{Caveat.} In clear weather the proposed model (0.925) is nominally ahead of the caution-aware model (0.890), a gap that is not significant ($p=0.092$). The energy saving is an estimate and the caution value is given to the model, not sensed by it.
\end{implicationbox}
